% Background additions — paste after Section 2.8 in chapters/15-background.tex.
% Requires the candidate references (Notes/overleaf/candidate_references.bib
% contents) merged into thesis.bib first.
% Revised 2026-09-07 after independent review (citation accuracy + prose).

\section{State of the Art in Neuromechanical Simulation}\label{sec:bg_sota}

Neuromechanical simulation couples a musculoskeletal plant to a model of the
neural controller that drives it. One line of work treats reflex physiology as
the controller. Geyer and Herr \citep{geyer_musclereflex_2010} showed that
positive force, length, and velocity feedback around each muscle, organized with
simple interlimb coordination, is enough to produce self-stabilizing human-like
walking dynamics and realistic muscle activations without a central pattern
generator. Complementary neuromechanical models of the cat hindlimb combine
central pattern generation with afferent control, as discussed in
Section~\ref{sec:bg_neural} \citep{markin_afferent_2010, shevtsova_modeling_2016}.

A second strand replaces hand-designed circuitry with machine learning. Song et
al. \citep{song_deep_2021} surveyed this direction and reported the Learn to
Move competitions, in which deep reinforcement learning policies were trained on
OpenSim musculoskeletal models and produced walking and running with realistic
kinematics. The plant in such work is mature, but the controller is an opaque
function approximation: its internal variables do not correspond to identifiable
physiological circuits, so the lesion, deletion, and stimulation experiments
that make neuromechanical models scientifically useful lose their meaning.
Reviews of predictive musculoskeletal simulation likewise discuss the persistent
difficulty of specifying and validating the assumed motor control
\citep{degroote_perspective_2021}.

This dissertation belongs to a third approach: synthetic nervous systems,
controllers built from explicit, conductance-based neurons mapped to identified
physiological circuits and simulated at speeds compatible with real-time
robotics. The program attempted here combines an anatomically grounded
controller, a biomechanically validated pneumatic plant
(Chapters~\ref{ch:methods}--\ref{ch:results}), and an eventual physical robot.
Simulation speed makes the combination practical. SNS-Toolbox couples directly
to physics simulators such as MuJoCo and sustains real-time or faster network
evaluation from workstations down to embedded boards
\citep{caggiano_myosuite_2022, nourse_sns_2023}. A single overnight run can
sweep thousands of parameter combinations or execute hundreds of simulated gait
cycles per condition. This throughput makes lesion and stimulation studies
statistically practicable on the model, at a scale that human-subject and
animal protocols rarely reach.

\section{Testing Neuromechanical Hypotheses on Robots: The Ethical Case}\label{sec:bg_ethics}

The experiments this research program is designed to enable cannot be performed
on healthy humans. Deleting a muscle from a limb, blocking a sensory pathway,
lesioning part of a rhythm generator, or applying tonic stimulation to spinal
circuits is invasive, irreversible, and confined in humans to clinical
populations under tight ethics review, without the within-subject control and
reversibility that neuromechanical hypotheses demand. The same experiments in
animals carry real moral cost. Their justification is governed by the Three Rs
framework of humane experimental technique: replacement of animal subjects with
non-animal methods, reduction of the numbers required, and refinement of the
remaining procedures \citep{russell_principles_1959, tannenbaum_russell_2015}. A
robot that implements a validated neuromechanical hypothesis is such a
replacement method. Every deletion, lesion, and stimulation experiment that runs
on the synthetic nervous system is an experiment that does not require a living
subject. Robot subjects can also be duplicated exactly, so the populations
across experimental conditions are statistically identical, which breeding
colonies cannot guarantee. Refinement follows directly: the failure modes of
these experiments damage hardware, not tissue.

Before a robot can stand in as a replacement subject, however, its validation
must meet an explicit standard. Webb's analyses of biorobotic methodology
\citep{webb_robots_2001, webb_robots_2002} argue that a robot model earns
biological relevance when it tests hypotheses about the target mechanism,
represents the structure of that mechanism accurately, and matches, and ideally
predicts, the target behavior, with relevance judged across several dimensions
rather than by surface resemblance. The validation chain of this dissertation
follows that standard. The actuator force model was validated against bench
data (Section~\ref{sec:results_force}), the joint torque model against
isometric measurements and human torque envelopes
(Section~\ref{sec:results_torque}), and the neuromechanical controller will be
verified against published tonic stimulation and deletion phenomena
\citep{selionov_tonic_2009, rybak_modelling_2006} before any locomotion tuning
is attempted. Only then can an experiment on the robot be read as evidence
about the hypothesized biological circuit.

These considerations return the argument to the program described in
Section~\ref{sec:project}. Experiments that are unethical, impractical, or
statistically out of reach in humans and animals become routine on the robot.
The resulting insights feed back into biology, into the design of the next
robot iteration, and into prosthetic and orthotic devices, which face the same
underlying problem of controlling muscle-like actuators for human users. The
ethical and scientific cases thus coincide: the robot is the instrument that
makes these experiments possible.
